\RequirePackage[!use=preview,!use=efpmca,!use=usebw]{spdef}
\documentclass{article}
\usepackage[fleqn]{amsmath}
\usepackage[
    web={centertitlepage,designv,forcolorpaper,
         latextoc,extended},
    eforms,aebxmp
]{aeb_pro}
\usepackage{ran_toks}

\useThisSeed{227819015}
%\useThisSeed{1441984427}
%\useLastAsSeed
%\useThisSeed{2001383846}

%2001383846 % initializing seed value
%1629639958 % last random number used

\let\pkg\textsf
\let\env\texttt
\let\opt\textsf

%\rtdebugtrue
%\ranToksOn
%\ranToksOff


%\usepackage{myriadpro}
\usepackage[altbullet]{lucidbry}

%\usepackage{makeidx}
%\makeindex
\usepackage{acroman}
\usepackage[active]{srcltx}

\makeatletter
\def\eq@fititin#1{\noindent\unskip\nobreak\hfill\penalty50
    \hskip2em\hbox{}\nobreak\hfill#1}
\def\fitit{\eq@fititin{\exrtnlabelformat}}
\@mparswitchfalse\reversemarginpar
\renewcommand{\paragraph}
    {\@startsection{paragraph}{4}{0pt}{6pt}{-3pt}
    {\normalfont\normalsize\bfseries}}
\renewcommand{\subparagraph}
    {\@startsection{subparagraph}{5}{\parindent}{6pt}{-3pt}%
    {\normalfont\normalsize\bfseries}}
\makeatother

\def\anglemeta#1{$\langle\textit{\texttt{#1}}\rangle$}
\let\ameta\anglemeta
\def\meta#1{\textit{\texttt{#1}}}
\def\darg#1{\texttt{\{#1\}}}
\def\takeMeasure{\bgroup\obeyspaces\takeMeasurei}
\def\takeMeasurei#1{\global\setbox\webtempboxi\hbox{\ttfamily#1}\egroup}
\def\bxSize{\wd\webtempboxi+2\fboxsep+2\fboxrule}

\newcommand{\FmtMP}[2][0pt]{\mbox{}\marginpar{%
    \raisebox{.5\baselineskip+#1}{%
    \expandafter\parbox\expandafter[\parboxValign]%
        {\marginparwidth}{\aebbkFmtMp#2}}}}
\def\aebbkFmtMp{\kern0pt\itshape\small
    \ifusebw\color{gray}\else\color{blue}\fi
    \raggedleft\hspace{0pt}}\newcommand{\BlogArticle}{%
    \makebox[0pt][l]{\hspace{-1pt}\color{blue}\Pisymbol{webd}{254}%
    }\raisebox{.5pt}{\ifusebw\color{black}\else
    \color{red}\fi\ding{045}}}
\def\parboxValign{t}


\usepackage[active]{srcltx}

\let\amtIndent\leftmargini
\def\SUB#1{${}_{\text{#1}}$}

\newdimen\aebdimen \aebdimen\topsep
\newcommand\bVerb[1][]{\begingroup#1\vskip\aebdimen\parindent0pt}%
\def\eVerb{\vskip\aebdimen\endgroup\noindent}


\urlstyle{rm}

\let\tops\texorpdfstring

\DeclareDocInfo
{
    university={\AcroTeX.Net},
    title={\tops{The \textsf{ran\_toks}}{The manual for the ran\_toks}
        Package\tops{\\[6pt]\large}{: }
        Randomizing the order of tokens},
    author={D. P. Story},
    email={dpstory@acrotex.net},
    subject=Documentation for the ran\_toks package,
    talksite={\url{www.acrotex.net}},
    version={1.4, 2021/06/06},
    Keywords={LaTeX, PDF, random, tokens, JavaScript, Adobe Acrobat},
    copyrightStatus=True,
    copyrightNotice={Copyright (C) \the\year, D. P. Story},
    copyrightInfoURL={http://www.acrotex.net}
}
\DeclareInitView{windowoptions={showtitle}}


\def\dps{$\hbox{$\mathfrak D$\kern-.3em\hbox{$\mathfrak P$}%
       \kern-.6em \hbox{$\mathcal S$}}$}

\universityLayout{fontsize=Large}
\titleLayout{fontsize=LARGE}
\authorLayout{fontsize=Large}
\tocLayout{fontsize=Large,color=aeb}
\sectionLayout{indent=-62.5pt,fontsize=large,color=aeb}
\subsectionLayout{indent=-31.25pt,color=aeb}
\subsubsectionLayout{indent=0pt,color=aeb}
\subsubDefaultDing{\tops{$\bullet$}{\textrm\textbullet}}

\def\exSrc{\makebox[0pt][r]{\large{\Pisymbol{webd}{157}}\enspace}}

%\pagestyle{empty}
%\parindent0pt\parskip\medskipamount

\chngDocObjectTo{\newDO}{doc}
\begin{docassembly}
var titleOfManual="The ran_toks Package";
var manualfilename="Manual_BG_Print_rt.pdf";
var manualtemplate="Manual_BG_Brown.pdf"; // Blue, Green, Brown
var _pathToBlank="C:/Users/Public/Documents/ManualBGs/"+manualtemplate;
var doc;
var buildIt=false;
if ( buildIt ) {
    console.println("Creating new " + manualfilename + " file.");
    doc = \appopenDoc({cPath: _pathToBlank, bHidden: true});
    var _path=this.path;
    var pos=_path.lastIndexOf("/");
    _path=_path.substring(0,pos)+"/"+manualfilename;
    \docSaveAs\newDO ({ cPath: _path });
    doc.closeDoc();
    doc = \appopenDoc({cPath: manualfilename, oDoc:this, bHidden: true});
    f=doc.getField("ManualTitle");
    f.value=titleOfManual;
    doc.flattenPages();
    \docSaveAs\newDO({ cPath: manualfilename });
    doc.closeDoc();
} else {
    console.println("Using the current "+manualfilename+" file.");
}
var _path=this.path;
var pos=_path.lastIndexOf("/");
_path=_path.substring(0,pos)+"/"+manualfilename;
\addWatermarkFromFile({
    bOnTop:false,
    bOnPrint:false,
    cDIPath:_path
});
\executeSave();
\end{docassembly}

%\definePath\bgPath{"C:/Users/Public/Documents/%
%    ManualBGs/Manual_BG_Print_AeB.pdf"}
%\begin{docassembly}
%\addWatermarkFromFile({%
%    bOnTop: false,
%    cDIPath: \bgPath
%})
%\executeSave()
%\end{docassembly}

\begin{document}

\maketitle

\selectColors{linkColor=black}
\tableofcontents
\selectColors{linkColor=webgreen}

\section{Introduction}

This is a short package for randomizing the order of tokens. The package
is long overdue; users of \textsf{\textbf{AeB}} and of \pkg{eqexam} have long asked for a way to
randomize the order of the problems in a test or quiz, or anything for
that matter.

\newtopic\noindent\exSrc The \texttt{examples} folder contains nine demonstration files:
\begin{enumerate}
    \item \texttt{ran-toks.tex} reproduces the sample code of this manual.
    \item \texttt{rt-tst-eqe.tex} shows how to use \pkg{ran\_toks} to
        randomize the \emph{questions} of an exam document created by the
        \pkg{eqexam} package.
    \item \texttt{rt-tst-qz.tex} shows how to randomize choices of a
        multiple choice field in a \env{quiz} environment of the
        \pkg{exerquiz} package, when the choices contain verbatim text.
    \item \texttt{rt-cb.tex} is a demonstration of how to use the
        \pkg{collectbox} package to place verbatim content into a token list
        (see \cs{ranToks}, \hyperref[s:ranToks]{Section~\ref*{s:ranToks}}).
\xdef\svEnumi{\the\value{enumi}}
\end{enumerate}
The next five files concern the DB application and require the \opt{usedbapp} option,
they are discussed in \hyperref[s:DBApp]{Section~\ref*{s:DBApp}}.
\begin{enumerate}\setcounter{enumi}{\svEnumi}
    \item \texttt{mc-db.tex} is an \pkg{eqexam} file, draws from the
        database files \texttt{db1.tex}, \texttt{db2.tex},
        \texttt{db3.tex}, and \texttt{db4.tex}, to construct the questions
        of the exam. The questions are drawn at random from the DB files. Refer
        to \hyperref[s:DBApp]{Section~\ref*{s:DBApp}} for a few more details.
    \item \texttt{mc-dbu.tex} is an \pkg{eqexam} file that draws from the
        database test file, draws from the database files \texttt{db1.tex},
        \texttt{db2.tex}, \texttt{db3.tex}, and \texttt{db4.tex}, to
        construct the questions of the exam. However, \texttt{mc-dbu.tex}
        differs from \texttt{mc-db.tex} for it selects unique questions
        cross the four versions of the document. Refer to
        \hyperref[s:DBApp]{Section~\ref*{s:DBApp}} for a few more details.
    \item \texttt{viewDB.tex} A demo file of how to view the DB files for
        the purpose of reviewing the questions, modifying them, or adding
        to them.
    \item \texttt{mc-dbu-ctrld.tex} is a variation on \texttt{mc-dbu.tex},
        it is not meant to be compiled directly, but by the file
        \texttt{ctrl-build.tex}.
    \item \texttt{ctrl-build.tex} is a simple TEX file that controls the
        compilation of a targeted file (\texttt{mc-dbu-ctrld.tex} in this
        example) over the various versions supported by the targeted file.
        Try it, you'll like it. Uses the \pkg{shellesc} package to control
        the process.

\end{enumerate}

\section{The Preamble and Package Options}

The preamble for this package is
\bVerb\takeMeasure{\string\usepackage[\ameta{options}]\darg{ran\_toks}}%
\begin{minipage}{\bxSize}\kern0pt
\begin{Verbatim}[frame=single,commandchars=!()]
\usepackage[!ameta(options)]{ran_toks}
\end{Verbatim}
\end{minipage}\eVerb
The one option of \pkg{ran\_toks} is \opt{usedbapp}\FmtMP{\opt{usedbapp} option}, which brings in
specialized code to support a database test application. This option is
discussed in \hyperref[s:DBApp]{Section~\ref*{s:DBApp}}, beginning on page~\pageref*{s:DBApp}.

The requirements for \textsf{ran\_toks} are the \textsf{verbatim} package
(part of the standard {\LaTeX} distribution, and the macro file
\texttt{random.tex} by Donald Arseneau.


\section{The main commands and environments}\label{rtmain}

There are two styles for defining a series of tokens to be randomized,
using either the \cs{ranToks} command or the \cs{bRTVToks}/\cs{eRTVToks}
pair. Each of these is discussed in the next two subsections.

\subsection{The \tops{\protect\cs{ranToks}}{\CMD{ranToks}} command}\label{s:ranToks}

The \cs{ranToks} command was the original concept; declare a series of
tokens to be randomized.
\bVerb\takeMeasure{\string\ranToks\darg{\ameta{name}}\{\%\enspace}%
\begin{minipage}{\bxSize}\kern0pt
\begin{Verbatim}[frame=single,commandchars=!()]
\ranToks{!ameta(name)}{%
!quad{!ameta(token!SUB(1))}
!quad{!ameta(token!SUB(2))}
!quad...
!quad{!ameta(token!SUB(n))}
}
\end{Verbatim}
\end{minipage}\eVerb
were \ameta{token\SUB{k}} is any non-verbatim content;\footnote{However, for workarounds, see
\texttt{rt-cb.tex} and \texttt{rt-tst-qz.tex.}} each token is enclosed in braces
(\darg{}), this is required. The \ameta{name} parameter is required, and
must be unique for the document; it is used to build the names of internal
macros. Of course several such \cs{ranToks} can be used in the document,
either in the preamble or in the body of the document. Multiple
\cs{ranToks} commands must have a different \ameta{name} parameter.

\emph{After} a \cs{ranToks} command has been executed, the number of tokens
counted is accessible through the \cs{nToksFor} command,
\bVerb\takeMeasure{\string\nToksFor\darg{\ameta{name}}}%
\begin{minipage}{\bxSize}\kern0pt
\begin{Verbatim}[frame=single,commandchars=!()]
\nToksFor{!ameta(name)}
\end{Verbatim}
\end{minipage}\eVerb
The one argument is \ameta{name}, and will expand to the total number
of tokens listing as argument in the \cs{ranToks} command by the same
name.

The \cs{ranToks} command does not display the randomized tokens, for that
the command \cs{useRanTok} is used.
\bVerb\takeMeasure{\string\useRTName\darg{\ameta{name}}}%
\begin{minipage}{\bxSize}
\begin{Verbatim}[frame=single,commandchars=!()]
\useRanTok{!ameta(num)}
\useRTName{!ameta(name)}
\end{Verbatim}
\end{minipage}\eVerb
The argument of \cs{useRanTok} is a positive integer between 1 and
\cs{nToksFor\darg{\ameta{name}}}, the number of tokens declared by
\cs{ranToks}, inclusive. There is no space created following the
\cs{useRanTok} command, so if these are to be used ``inline'', enclose them
in braces (\darg{}), for example, \darg{\cs{useRanTok\darg{1}}}. The use of
\cs{useRTName} is optional unless the listing of the \cs{useRanTok} commands
is separated from the \cs{ranToks} command that defined them by another
\cs{ranToks} command of a different name. That should be clear!

Consider this example.

\begin{Verbatim}[xleftmargin=\amtIndent]
\ranToks{myPals}{%
    {Jim}{Richard}{Don}
    {Alex}{Tom}{J\"{u}rgen}
}
\end{Verbatim}
\ranToks{myPals}{%
    {Jim}{Richard}{Don}
    {Alex}{Tom}{J\"{u}rgen}
}
I have {\nToksFor{myPals}} pals, they are \useRanTok{1}, \useRanTok{2},
\useRanTok{3}, \useRanTok{4}, {\useRanTok{5}} and \useRanTok{6}. (Listed
in the order of best friend to least best friend.)
The verbatim listing is,
\begin{Verbatim}[xleftmargin=\amtIndent]
I have {\nToksFor{myPals}} pals, they are \useRanTok{1},
\useRanTok{2}, \useRanTok{3}, \useRanTok{4}, {\useRanTok{5}}
and \useRanTok{6}.
\end{Verbatim}
Notice that \cs{useRanToks} are not enclosed in braces for 1--4 because
they are each followed by a comma; the fifth token, \texttt{\darg{\cs{useRanTok\darg{5}}}}, is
enclosed in braces to generate a space following the insertion of the text.

Repeating the sentence yields, ``I have {\nToksFor{myPals}} pals, they are
\useRanTok{1}, \useRanTok{2}, \useRanTok{3}, \useRanTok{4},
{\useRanTok{5}} and \useRanTok{6}'', which is the exact same random order. To
obtain a different order, re-execute the \cs{ranToks} command with the same
arguments.\footnote{{\ttfamily\string\ranToks\darg{myPals}\darg{\darg{Jim}\darg{Richard}\darg{Don}\darg{Alex}\darg{Tom}\darg{J\string\"\darg{u}rgen}}} in this example.} Doing just that, \ranToks{myPals}{{Jim}{Richard}{Don}
{Alex}{Tom}{J\"{u}rgen}}we obtain, ``I have {\nToksFor{myPals}} pals, they
are \useRanTok{1}, \useRanTok{2}, \useRanTok{3}, \useRanTok{4},
{\useRanTok{5}} and \useRanTok{6}.'' A new order?
An alternative to re-executing \cs{ranToks} is to use the \cs{reorderRanToks} command:
\bVerb\takeMeasure{\string\reorderRanToks\darg{\ameta{name}}}%
\begin{minipage}{\bxSize}\kern0pt
\begin{Verbatim}[frame=single,commandchars=!()]
\reorderRanToks{!ameta(name)}
\end{Verbatim}
\end{minipage}\eVerb Now, executing \verb~\reorderRanToks{myPals}~ and compiling the sentence again yields,
``I~have {\nToksFor{myPals}} pals, they are
\useRanTok{1}, \useRanTok{2}, \useRanTok{3}, \useRanTok{4},
{\useRanTok{5}} and \useRanTok{6}.'' For most applications, re-randomizing
the same token list in the same document is not very likely something you
need to do.

The \cs{reorderRanToks\darg{\ameta{name}}} command rearranges the list of tokens
associated with \ameta{name}, which may not be what you want; the
\cs{copyRanToks} command, on the other hand, makes a (randomized) copy of its
first required argument \ameta{name\SUB1} and saves it as \ameta{name\SUB2}, without
effecting the order of \ameta{name\SUB1}.
\bVerb\takeMeasure{\string\copyRanToks\darg{\ameta{name\SUB1}}\darg{\ameta{name\SUB2}}}%
\begin{minipage}{\bxSize}\kern0pt
\begin{Verbatim}[frame=single,commandchars=!()]
\copyRanToks{!ameta(name!SUB1)}{!ameta(name!SUB2)}
\end{Verbatim}
\end{minipage}\eVerb Thus, if \verb!\copyRanToks{myPals}{myPals1}! is executed, the token list name
\texttt{myPals1} contains the names of my pals in another randomized order,
while maintaining the same order of \texttt{myPals}.


\newtopic
My original application for this, the one that motivated writing this
package at long last, was the need to arrange several form buttons
randomly on the page. My point is that the listing given in the argument
of \cs{ranToks} can pretty much be anything that is allowed to be an
argument of a macro; this would exclude verbatim text created by \cs{verb}
and verbatim environments.

\subsection{The
\tops{\protect\cs{bRTVToks}/\protect\cs{eRTVToks}}
    {\CMD{bRTVToks}/\CMD{eRTVToks}} pair of commands}\label{ss:RTVToks}

Sometimes the content to be randomized is quite large or contains verbatim
text. For this, it may be more convenient to use the
\cs{bRTVToks}/\penalty0\cs{eRTVToks} command pair. The syntax is
\bVerb\takeMeasure{\string\bRTVToks\darg{\ameta{name}}}%
%\edef\WIDTH{\the\wd\webtempboxi}%
\def\1{\rlap{\hspace*{\linewidth}\texttt{\%{\sffamily{ \raisebox{1pt}{<--}Begin token listing}}}}}%
\def\2{\rlap{\hspace*{\linewidth}\texttt{\%{\sffamily{ \raisebox{1pt}{<--}End token listing}}}}}%
%\takeMeasure{\string\bRTVToks\darg{\ameta{name}}\quad\%{\sffamily{ <-Begin token listing}}}%
\begin{minipage}{\bxSize}\kern0pt
\begin{Verbatim}[frame=single,commandchars=!()]
!1\bRTVToks{!ameta(name)}
\begin{rtVW}
!quad!ameta(content!SUB(1))
\end{rtVW}
...
...
\begin{rtVW}
!quad!ameta(content!SUB(n))
\end{rtVW}
!2\eRTVToks
\end{Verbatim}
\end{minipage}\eVerb
%\bVerb\takeMeasure{\string\bRTVToks\darg{\ameta{name}}\quad}%
%\edef\WIDTH{\the\wd\webtempboxi}%
%\def\1{\rlap{\hspace*{\WIDTH}\texttt{\% <-{\sffamily{ End token listing}}}}}%
%\takeMeasure{\string\bRTVToks\darg{\ameta{name}}\quad\% <-{\sffamily{ Begin token listing}}}%
%\begin{minipage}{\bxSize}\kern0pt
%\begin{Verbatim}[frame=single,commandchars=!()]
%\bRTVToks{!ameta(name)}!quad% <-!sffamily( Begin token listing)
%\begin{rtVW}
%!quad!ameta(content!SUB(1))
%\end{rtVW}
%...
%...
%\begin{rtVW}
%!quad!ameta(content!SUB(n))
%\end{rtVW}
%!1\eRTVToks
%\end{Verbatim}
%\end{minipage}\eVerb
The \cs{bRTVToks}\texttt{\darg{\ameta{name}}} command begins the (pseudo)
environment and is ended by \cs{eRTVToks}. Between these two are a series of
\texttt{rtVW} (random toks verbatim write) environments. When the document is
compiled, the contents (\ameta{content\SUB{i}}) of each of these environments
are written to the computer hard drive and saved under a different name
(based on the parameter \ameta{name}). Later, using the \cs{useRanTok}
commands, they are input back into the document in a random order.

The \env{rtVW} environment also writes the command
\cs{RTVWHook}\marginpar{\small\itshape\flushright\cs{RTVWHook}} to the top of
the file. Its initial value is \cs{relax}. It can be redefined with
\cs{rtVWHook\darg{\ameta{arg}}}, which expands to
\verb|\def\RTVWHook{|\ameta{arg}\verb|}|.

The use of \cs{useRTName} and \cs{useRanTok} were explained and illustrated
in the previous section. Let's go to the examples,
\begin{Verbatim}[xleftmargin=\amtIndent]
\bRTVToks{myThoughts}
\begin{rtVW}
\begin{minipage}[t]{.67\linewidth}
Roses are red and violets are blue,
I've forgotten the rest, have you too?
\end{minipage}
\end{rtVW}
\begin{rtVW}
\begin{minipage}[t]{.67\linewidth}
I gave up saying bad things like
\verb!$#%%%^*%^&#$@#! when I was just a teenager.
\end{minipage}
\end{rtVW}
\begin{rtVW}
\begin{minipage}[t]{.67\linewidth}
I am a good guy, pass it on! The code for this last sentence is,
\begin{verbatim}
%#$% I am a good guy, pass it on! ^&*&^*
\end{verbatim}
How did that other stuff get in there?
\end{minipage}
\end{rtVW}
\eRTVToks
\end{Verbatim}
OK, now, let's display these three in random order. Here we place them in
an \texttt{enumerate} environment.

\bRTVToks{myThoughts}%
\begin{rtVW}
\begin{minipage}[t]{.67\linewidth}
Roses are red and violets are blue,
I've forgotten the rest, have you too?
\end{minipage}
\end{rtVW}
\begin{rtVW}
\begin{minipage}[t]{.67\linewidth}
I gave up saying bad things like
\verb!$#%%%^*%^&#$@#! when I was just a teenager.
\end{minipage}
\end{rtVW}
\begin{rtVW}
\begin{minipage}[t]{.67\linewidth}
I am a good guy, pass it on! The code for this last sentence is,
\begin{verbatim}
%#$% I am a good guy, pass it on! ^&*&^*
\end{verbatim}
How did that other stuff get in there?
\end{minipage}
\end{rtVW}
\eRTVToks
\begin{enumerate}
    \item \useRanTok{1}
    \item \useRanTok{2}
    \item \useRanTok{3}
\end{enumerate}
The verbatim listing of the example above is
\begin{Verbatim}[xleftmargin=\amtIndent]
\begin{enumerate}
    \item \useRanTok{1}
    \item \useRanTok{2}
    \item \useRanTok{3}
\end{enumerate}
\end{Verbatim}
The \cs{reorderRanToks} works for lists created by the \cs{bRTVToks}/\penalty0\cs{bRTVToks} construct.
If we say \cs{reorderRanToks\darg{myThoughts}} and reissue the above list, we obtain,
\begin{enumerate}\ranToksOn\reorderRanToks{myThoughts}
    \displayListRandomly[\item]{myThoughts}
\end{enumerate}
The command \cs{copyRanToks} works for list created by
\cs{bRTVToks}/\penalty0\cs{bRTVToks} as well.

\paragraph*{On the \cs{displayListRandomly} command.}\label{para:DLR}
In the enumerate example immediately above, the items in the list are
explicitly listed as \cs{item \cs{useRanTok\darg{1}}} and so on; an
alternate approach is to use the command \cs{displayListRandomly}, like so,
\begin{Verbatim}[xleftmargin=\amtIndent]
\begin{enumerate}
    \displayListRandomly[\item]{myThoughts}
\end{enumerate}
\end{Verbatim}
The full syntax for \cs{displayListRandomly} is displayed next.
\bVerb\takeMeasure{\string\displayListRandomly[\ameta{prior}][\ameta{post}]\darg{\ameta{name}}}%
\begin{minipage}{\bxSize}\kern0pt
\begin{Verbatim}[frame=single,commandchars=!()]
\displayListRandomly[!ameta(prior)][!ameta(post)]{!ameta(name)}
\end{Verbatim}
\end{minipage}
\eVerb The action of \cs{displayListRandomly} is to expand all tokens that
are listed in the \ameta{name} token list, each entry is displayed as
\ameta{prior}\cs{useRanTok\darg{i}}\ameta{post}, where \texttt{i}
goes from~1 to \cs{nToksFor\darg{\ameta{name}}}. In the example above,
\ameta{prior} is \cs{item}, but normally, its default is empty. The defaults
for \ameta{prior} and \ameta{post} are both empty.

\subparagraph*{The optional arguments.} When only one optional argument is present,
if is interpreted as \ameta{prior}. To obtain a \ameta{post} with no \ameta{prior}
use the syntax,
\begin{quote}
\cs{displayListRandomly[][\ameta{post}]\darg{\ameta{name}}}
\end{quote}
Within \emph{each optional argument}, the four commands \cs{i},
\cs{first}, \cs{last}, and \cs{lessone} are (locally) defined. The \cs{i} command is the index
counter of the token currently being typeset; \cs{first} is the index of the
first item; \cs{last} is the index of the last item; and \cs{lessone} is one
less than \cs{last}. The two optional arguments and the four commands may use
to perform logic on the token as it is being typeset. For example:
\begin{Verbatim}[xleftmargin=\amtIndent,fontsize=\small]
List of pals: \displayListRandomly
    [\ifnum\i=\last and \fi]
    [\ifnum\i=\last.\else, \fi]{myPals}
\end{Verbatim}
yields,
\begin{quote}
List of pals: \displayListRandomly
    [\ifnum\i=\last and \fi]
    [\ifnum\i=\last.\else, \fi]{myPals}
\end{quote}
The optional arguments are wrapped to the next line to keep them within the margins, cool.

The example above shows the list of my pals with an Oxford comma. How would
you modify the optional argument to get the same listing without the Oxford
comma? (\displayListRandomly[\ifnum\i=\last and \fi][\ifnum\i=\last.\else\ifnum\i=\lessone\relax\space\else, \fi\fi]{myPals})
Hint: a solution involves the other command \cs{lessone}.

\subsubsection{Nested
\tops{\protect\cs{bRTVToks}/\protect\cs{eRTVToks}}
    {\CMD{bRTVToks}/\CMD{eRTVToks}} command pairs}\label{ss:NestedRTVToks}

These is at least one example of using nested
\cs{bRTVToks}/\allowbreak\cs{eRTVToks}. When nested
\cs{bRTVToks}/\allowbreak\cs{eRTVToks} command pairs, use the \env{rtVWi}
environment instead of the \env{rtVW} environment. The nested
\cs{bRTVToks}/\allowbreak\cs{eRTVToks} pair is placed\FmtMP{placement} within
a \env{rtWV} environment; in this way the contents of that \env{rtWV}
environment, itself can be randomized. The \cs{displayListRandomly} (or
\cs{useRanTok}) command is used to list out the nested items. See the next
page for an example.

\newpage

\bRTVToks{Depth0}
  \begin{rtVW}
    Depth0: Item 1
\end{rtVW}
  \begin{rtVW}
    Depth0: Item 2:
    \bRTVToks{Depth1}
      \begin{rtVWi}
        Depth1: Item 1
\end{rtVWi}
      \begin{rtVWi}
        Depth1: Item 2
        \bRTVToks{Depth2}
          \begin{rtVWii}
              Depth2: Item 1
\end{rtVWii}
          \begin{rtVWii}
              Depth2: Item 2
\end{rtVWii}
          \begin{rtVWii}
              Depth2: Item 3
\end{rtVWii}
        \eRTVToks
\begin{itemize}
        \displayListRandomly[\item]{Depth2}
\end{itemize}
\end{rtVWi}
      \begin{rtVWi}
        Depth1: Item 3
\end{rtVWi}
\eRTVToks
\begin{itemize}
\displayListRandomly[\item]{Depth1}
\end{itemize}
\end{rtVW}
  \begin{rtVW}
      Depth0: Item 3
\end{rtVW}
  \begin{rtVW}
      Depth0: Item 4
\end{rtVW}
  \begin{rtVW}
      Depth0: Item 5
\end{rtVW}
\eRTVToks
   \begin{enumerate}\small
\displayListRandomly[\item]{Depth0}
\end{enumerate}
The verbatim listing for this example follows.
\begin{Verbatim}[commandchars=!(),fontsize=\small]
\bRTVToks{Depth0}
  \begin{rtVW}
    Depth0: Item 1
\end{rtVW}
  \begin{rtVW}
    Depth0: Item 2:
    \bRTVToks{Depth1} % !textsf(begin Depth2 within an rtVW env.)
      \begin{rtVWi}
        Depth1: Item 1
\end{rtVWi}
      \begin{rtVWi}
        Depth1: Item 2
        \bRTVToks{Depth2} % !textsf(begin Depth1 within an rtVWi env.)
          \begin{rtVWii}
              Depth2: Item 1
\end{rtVWii}
          \begin{rtVWii}
              Depth2: Item 2
\end{rtVWii}
          \begin{rtVWii}
              Depth2: Item 3
\end{rtVWii}
        \eRTVToks
\begin{itemize}
        \displayListRandomly[\item]{Depth2} % !textsf(display Depth2 toks)
\end{itemize}
\end{rtVWi}
      \begin{rtVWi}
        Depth1: Item 3
\end{rtVWi}
\eRTVToks
\begin{itemize}
\displayListRandomly[\item]{Depth1} % !textsf(display Depth1 toks)
\end{itemize}
\end{rtVW}
  \begin{rtVW}
      Depth0: Item 3
\end{rtVW}
  \begin{rtVW}
      Depth0: Item 4
\end{rtVW}
  \begin{rtVW}
      Depth0: Item 5
\end{rtVW}
\eRTVToks
\begin{enumerate}
\displayListRandomly[\item]{Depth0} % !textsf(display Depth0 toks)
\end{enumerate}
\end{Verbatim}
Oops. Did I forget to mention that \pkg{ran\_toks} supports nested to a depth of two.

Some authors like to indent nested things, but to avoid spurious spaces
appearing, \cs{end\darg{rtVW}} (and \cs{end\darg{rtVWi}}) should be placed in
the far left margin, as shown above. Recall the \env{rtVW} is a verbatim
environment.

The above example is reproduced in the \texttt{ran\_toks.tex} sample file found in the
\texttt{examples} folder. Also found in that folder is \texttt{nested-matching.tex}, a \pkg{exerquiz}
quiz that motivated creating nested \cs{bRTVToks}/\allowbreak\cs{eRTVToks} command pairs.

\section{Additional arguments and commands}\label{AddCmds}

The syntax given earlier for \cs{useRanTok} was not completely specified.
It is,
\bVerb\takeMeasure{\string\useRanTok[\ameta{name}]\darg{\ameta{num}}}%
\begin{minipage}{\bxSize}\kern0pt
\begin{Verbatim}[frame=single,commandchars=!()]
\useRanTok[!ameta(name)]{!ameta(num)}
\end{Verbatim}
\end{minipage}\eVerb
The optional first parameter specifies the \ameta{name} of the list from
which to draw a random token; \ameta{num} is the number of the
token in the range of 1 and \cs{nToksFor\darg{\ameta{name}}},
inclusive. The optional argument is useful in special circumstances when
you want to mix two random lists together.

\newtopic\noindent To illustrate: \useRanTok[myPals]{1}, \useRanTok[myThoughts]{2}

\newtopic\noindent The verbatim listing is
\begin{Verbatim}[xleftmargin=\amtIndent]
To illustrate: \useRanTok[myPals]{1}, \useRanTok[myThoughts]{2}
\end{Verbatim}
The typeset version looks a little strange, but recall, the text of
\texttt{myThoughts} were each put in a \texttt{minipage} of width \texttt{.67\cs{linewidth}}.
Without the \texttt{minipage}, the text would wrap around normally.

\paragraph*{Accessing the original order.} The original order of the list of tokens is not lost, you can retrieve
them using the command \cs{rtTokByNum},
\bVerb\takeMeasure{\string\rtTokByNum[\ameta{name}]\darg{\ameta{num}}}%
\begin{minipage}{\bxSize}\kern0pt
\begin{Verbatim}[frame=single,commandchars=!()]
\rtTokByNum[!ameta(name)]{!ameta(num)}
\end{Verbatim}
\end{minipage}\eVerb
This command expands to the token declared in the list named \ameta{name}
that appears at the \ameta{num} place in the list. (Rather awkwardly written.)
For example, my really best pals are {\rtTokByNum[myPals]{3}} and
\rtTokByNum[myPals]{4}, but don't tell them. The listing is,
\begin{Verbatim}[xleftmargin=\amtIndent]
For example, my really best pals are {\rtTokByNum[myPals]{3}}
and \rtTokByNum[myPals]{4}, but don't tell them.
\end{Verbatim}
In some sense, \cs{rtTokByNum[\ameta{name}]} acts like a simple array, the
length of which is \cs{nToksFor\{\ameta{name}\}}, and whose $k^{\text{th}}$
element is \cs{rtTokByNum[\ameta{name}]\{\ameta{k}\}}.

\paragraph*{Turning off randomization.} The randomization may be turned off
using \cs{ranToksOff} or turned back on with \cs{ranToksOn}.
\bVerb\takeMeasure{\string\ranToksOff\quad\string\ranToksOn}%
\begin{minipage}{\bxSize}\kern0pt
\begin{Verbatim}[frame=single,commandchars=!()]
\ranToksOff!quad\ranToksOn
\end{Verbatim}
\end{minipage}\eVerb
This can be done globally in the preamble for the whole of the document,
or in the body of the document just prior to either \cs{ranToks} or
\cs{bRTVToks}. For example,
\begin{Verbatim}[xleftmargin=\amtIndent]
\ranToksOff
\ranToks{integers}{ {1}{2}{3}{4} }
\ranToksOn
\end{Verbatim}
As a check, executing `$\cs{useRanTok\darg{3}} =  \cs{rtTokByNum\darg{3}} = 3 $' yields
`\ranToksOff\ranToks{integers}{ {1}{2}{3}{4}}\ranToksOn
$\useRanTok{3} = \rtTokByNum{3} = 3 $'? As anticipated.

To create a non-randomized list of tokens that already have been created (and randomized), use
\cs{copyRanToks}:
\begin{Verbatim}[xleftmargin=\amtIndent]
\ranToksOff\copyRanToks{myPals}{myOriginalPals}\ranToksOn
\end{Verbatim}
Then, using \cs{displayListRandomly} in a clever way,
\begin{Verbatim}[xleftmargin=\amtIndent]
\displayListRandomly[\ifnum\i=\last\space and \fi(\the\i)~]
    [\ifnum\i=\last.\else,\fi\space]{myOriginalPals}
\end{Verbatim}
we obtain: \ranToksOff\copyRanToks{myPals}{myOriginalPals}\ranToksOn
\displayListRandomly[\ifnum\i=\last\space and \fi(\the\i)~][\ifnum\i=\last.\else,\fi\space]{myOriginalPals}
The original list for \texttt{myPals} remains unchanged:
\displayListRandomly[\ifnum\i=\last\space and \fi(\the\i)~][\ifnum\i=\last.\else,\fi\space]{myPals}

The \cs{useRanTok} command---whether it operates on a randomized token list
or not---behaves similarly to an array. Thus, if we wanted to extract
the third entry of the non-randomized token list (array)
\texttt{myOriginalPals}, we do so by expanding the command
\verb!\useRanTok[myOriginalPals]{3}! to produce
\useRanTok[myOriginalPals]{3}.

\subparagraph*{Document preparation.}
The command \cs{ranToksOff} is probably best in the preamble to turn off
all randomization while the rest of the document is being composed.

\paragraph*{The \textsf{ran\_toks} auxiliary file.} The package writes to a file named
\cs{jobname\_rt.sav}, below represents two typical lines in this file.
\begin{Verbatim}[xleftmargin=\amtIndent]
1604051353 % initializing seed value
5747283528 % last random number used
\end{Verbatim}
The first line is the initializing seed value used for the last
compilation of the document; the second line is the last value of the
pseudo-random number generator used in the document.

Normally, the pseudo-random number generator provided by
\texttt{random.tex} produces a new initial seed value every minute. So if
you recompile again before another minute, you'll get the same initial
seed value.

\paragraph*{Controlling the initial seed value.}\label{para:ControlSeed} To obtain a new initial seed
value each time you compile, place \cs{useLastAsSeed} in the preamble.
\bVerb\takeMeasure{\string\useLastAsSeed}%
\begin{minipage}{\bxSize}\kern0pt
\begin{Verbatim}[frame=single]
\useLastAsSeed
\end{Verbatim}
\end{minipage}\eVerb
When the document is compiled, the initial seed value taken as the second
line in the \cs{jobname\_rt.sav} file, as seen in the above example.
With this command in the preamble, a new set of random numbers is
generated on each compile. If the file \cs{jobname\_rt.sav} does not
exist, the generator will be initialized by its usual method, using the time and date.

The command \cs{useThisSeed} allows you to reproduce a previous
pseudo-random sequence.
\bVerb\takeMeasure{\string\useThisSeed\darg{\ameta{init\_seed\_value}}}%
\begin{minipage}{\bxSize}\kern0pt
\begin{Verbatim}[frame=single,commandchars=!()]
\useThisSeed{!ameta(init_seed_value)}
\end{Verbatim}
\end{minipage}\eVerb
This command needs to be placed in the preamble. The value of
\ameta{init\_seed\_value} is an integer, normally taken from the
first line of the \cs{jobname\_rt.sav} file.

When creating tests (possibly using \textsf{eqexam}), the problems, or
contiguous collections of problems, can be randomly ordered using the
\cs{bRTVToks}/\penalty0\cs{eRTVToks} command pair paradigm. For example,
suppose there are two classes and you want a random order (some of) the
problems for each of the two classes. Proceed as follows:
\begin{enumerate}
\item Compile the document, open \cs{jobname\_rt.sav}, and copy the
    first line (in the above example, that would be
    \texttt{1604051353}).
\item Place \cs{useThisSeed\darg{1604051353}} in the preamble. Compiling
    will bring back the same pseudo-random sequence very time.
\item Comment this line out, and repeat the process (use
    \cs{useLastAsSeed} to generate new random sequences at each
    compile) until you get another distinct randomization, open
    \cs{jobname\_rt.sav}, and copy the first line again, say its \texttt{735794511}.
\item Place \cs{useThisSeed\darg{735794511}} in the preamble.
\item Label each
\begin{Verbatim}
%\useThisSeed{1604051353} % 11:00 class
%\useThisSeed{735794511}  % 12:30 class
\end{Verbatim}
To reproduce the random sequence for the class, just uncomment the random
seed used for that class.
\end{enumerate}
If you are using \textsf{eqexam}, the process can be automated as follows:
\begin{Verbatim}[xleftmargin=\amtIndent,commandchars=!()]
\vA{\useThisSeed{1604051353}} % 11:00 class
\vB{\useThisSeed{735794511}}  % 12:30 class
\end{Verbatim}
Again, this goes in the preamble.

\section{Commands to support a DB application}\label{s:DBApp}

The commands described in this section are only available with the \opt{usedbapp} option,
\begin{Verbatim}
\usepackage[usedbapp]{ran_toks}
\end{Verbatim}
This option was designed for an \pkg{eqexam} document.

\subsection{Basic functionality}\label{s:basicFunc}

The premise here is that you have an \pkg{eqexam} document (a test) and you have a series
of standard questions you ask students. Over the years, you have accumulated questions of a similar type
that you like to pose to your students. The questions of a similar type are placed in a DB test file. For example,
you have a file named \texttt{db1.tex} containing questions on a certain narrow topic. The format for this file
is,
\begin{Verbatim}[commandchars=!()]
%
% !textsf(Questions on some narrow topic)
%
\bRTVToks{DB1} %<-!textsf(DB !ameta(name))
\begin{rtVW}
% !textsf(an !pkg(eqexam) question is contained in this rtVW environment)
\end{rtVW}
\begin{rtVW}
% !textsf(an !pkg(eqexam) question is contained in this rtVW environment)
\end{rtVW}
...
\eRTVToks
\end{Verbatim}
Refer to the file \texttt{db1.tex}, \texttt{db2.tex},\dots, \texttt{db4.tex} for more specific
examples. The DB \ameta{name} must be unique among all the DB test files used.

The next step is to input your DB files. To do this, execute either of the commands
\cs{useTheseDBs} or \cs{useProbDBs} prior to the opening of an \env{exam} environment (\pkg{eqexam}),
or perhaps in the preamble. The syntax is,
\bVerb\takeMeasure{\string\useTheseDBs\darg{\ameta{db\SUB{1}},\ameta{db\SUB{2}},...,\ameta{db\SUB{n}}}}%
\setlength\eflength{\bxSize}%
\begin{minipage}{\bxSize}\def\1{\rlap{\sffamily\hskip\eflength(An alias for \cs{useTheseDBs})}}\kern0pt
\begin{Verbatim}[frame=single,commandchars=!()]
\useTheseDBs{!ameta(db!SUB(1)),!ameta(db!SUB(2)),...,!ameta(db!SUB(n))}
!1\useProbDBs{!ameta(db!SUB(1)),!ameta(db!SUB(2)),...,!ameta(db!SUB(n))}
\end{Verbatim}
\end{minipage}\eVerb The argument of \cs{useTheseDBs} is a comma-delimited
list of file names. Each file name contains a
\cs{bRTVToks}/\penalty0\cs{eRTVToks} construct. Within this pair are
\env{rtVW} environments, as described in
\hyperref[ss:RTVToks]{Section~\ref*{ss:RTVToks}}. The \cs{useTheseDBs}
command inputs the files listed in its comma-delimited argument; a warning is
emitted if one or more of the files are not found. The default extension is
\texttt{.tex}, \cs{useTheseDBs\darg{db1,db2}} inputs the files
\texttt{db1.tex} and \texttt{db2.tex}, if they exist, while
\cs{useTheseDBs\darg{db1.def,db2.db}} inputs the files \texttt{db1.def} and
\texttt{db2.db}, if they exist. The command \cs{useProbDBs} is an alias for
\cs{useTheseDBs}.

Once the database files have been input, in the body of the document you can extract
questions using \cs{useRanTok}; for example, \cs{useRanTok[DB1]\darg{1}} extracts the question
with an index of 1 from the database file DB1;\footnote{It is assumed that DB1 is the \ameta{name}
is the required argument of \cs{bRTVToks}, which is declared in the \texttt{db1.tex} file.} \cs{useRanTok[DB2]\darg{2}} extracts the question
with an index of 2 from the database file DB2, and so on.

Refer to the demonstration file \texttt{mc-db.tex} for an example.

\subsection{Striving for uniqueness of choices}

By default, when you expand \cs{useRanTok\darg{1}} and later you expand
\cs{useRanTok\darg{1}} again you get the exact same result;\footnote{Here, we
are assuming the \ameta{name} of the token list is the same for both
expansions of \cs{useRanTok\darg{1}}.} this is the normal behavior. However,
in the context of posing questions from a database of questions, repeating
the questions is not what is wanted necessarily. In the next two subsections,
we speak to the problem, first within a single-version document, and
secondly, within a multi-version document.


\subsubsection{For documents with a single version}

The issue of not wanting to repeat a question from a database comes up when
you want to pose several questions from a given database file. We illustrate
using the demo file \texttt{mc-db.tex} found in the \texttt{examples} folder.
In that file, the exam has two parts, as shown in
\hyperref[fig:2prts]{Figure~\ref*{fig:2prts}}.\footnote{The verbatim listing
here has been simplified.}
\begin{figure}[htb]\centering
\settowidth{\eflength}{\ttfamily\string\useRanTok[DB1]\darg{1}}%
\edef\x{\the\eflength}%
\setlength\eflength{\x+\widthof{\textsf{\quad\%An alternate approach}\texttt{ }}}%
\def\0{\rlap{\hskip\x\relax\quad\% \textsf{An alternate approach}}}%
\def\1{\rlap{\hskip\x\relax\quad\% \cs{begin\darg{Part2}}}}%
\def\2{\rlap{\hskip\x\relax\quad\% \cs{useRanTok[DB1]\darg{3}}}}%
\def\3{\rlap{\hskip\x\relax\quad\% \cs{useRanTok[DB1]\darg{4}}}}%
\def\4{\rlap{\hskip\x\relax\quad\% \cs{end\darg{Part2}}}}%
\fboxsep0pt
\begin{minipage}{\eflength}
\begin{Verbatim}[commandchars=!()]
\begin{Part1}
\useRanTok[DB1]{1}
\useRanTok[DB1]{2}
\end{Part1}

!0% !textsf( First approach)
!1\begin{Part2}
!2\useRanTok[DB1]{1}
!3\useRanTok[DB1]{2}
!4\end{Part2}
\end{Verbatim}
\end{minipage}
\caption{Simplified two-part exam}\label{fig:2prts}
\end{figure}
For \textsf{Part1}, we take two problems (at random) from \textsf{DB1}. In
\textsf{Part2}, we take two problems again from \textsf{DB1}; these two will
be the \emph{same} as the ones chosen from \textsf{Part1}. To get two
\emph{different} questions the natural approach is to write
\cs{useRanTok[DB1]\darg{3}} and \cs{useRanTok[DB1]\darg{4}}, as shown on the
right in \hyperref[fig:2prts]{Figure~\ref*{fig:2prts}}. Now, returning to the \textsf{First approach}, there is a way of forcing
\pkg{ran\_toks} to choose two different question even though the indexes used are the same.

To have different questions appear when you specify \verb~\useRanTok[DB1]{1}~
and \verb~\useRanTok[DB1]{2}~, first expand the command
\cs{uniqueXDBChoicesOn}, perhaps between parts or in the preamble. When
\cs{uniqueXDBChoicesOn} is expanded, \cs{useRanTok} tries to find an
``unused'' choice.
\begin{Verbatim}[xleftmargin=\parindent,commandchars=!()]
\uniqueXDBChoicesOn
\begin{document}
...
\begin{Part1}
\useRanTok[DB1]{1}
\useRanTok[DB1]{2}
\end{Part1}

%\uniqueXDBChoicesOn

\begin{Part2}
\useRanTok[DB1]{1}
\useRanTok[DB1]{2}
\end{Part2}
\end{Verbatim}
The command \cs{uniqueXDBChoicesOff} turns off the feature of striving to
find ``unused'' choices. Additional discussion on \cs{uniqueXDBChoicesOn} is
found in the next section; the command \cs{InputUsedIDs}, discussed in the
next section, is not needed when the \pkg{eqexam} document \emph{does not
have} multiple versions. You can experiment with striving to find unused
questions in the demo file \texttt{mc-db.tex}.

\subsubsection{For documents with multiple versions}

The scheme outlined in \hyperref[s:basicFunc]{Section~\ref*{s:basicFunc}} works well for an
\pkg{eqexam} document that only has one version in the source file. The way
\pkg{ran\_toks} works, it will not repeat random choices --- unless you
sample from a same DB file more times than there are problems in that file;
for example, suppose \texttt{db1.tex} has two questions in it, if you execute
\cs{useRanTok[DB1]\darg{1}}, \cs{useRanTok[DB1]\darg{2}}, and
\cs{useRanTok[DB1]\darg{3}}, then the problem selected by
\cs{useRanTok[DB1]\darg{3}} is the same as \cs{useRanTok[DB1]\darg{1}}. This
latter situation is not likely, is it?

One of the very powerful features of \pkg{eqexam} is that a single source
file can have multiple versions in it.
\begin{Verbatim}[xleftmargin=\parindent,commandchars=!()]
\examNum{1}
\numVersions{4}
\forVersion{a} %!textsf( a, b, c, d)
\vA{\useThisSeed{54356}}
\vB{\useThisSeed{577867}}
\vC{\useThisSeed{6746788}}
\vD{\useThisSeed{856785}}
\end{Verbatim}
The above shows how to set up a multi-version \pkg{eqexam} document, see the
\pkg{eqexam} documentation for more details. Here, as in the demo file
\texttt{mc-dbu.tex}, we declare 4 versions (a, b, c, and d).
\cs{forVersion\darg{a}} declares the next compile is for version a (or A,
either one). List also is a method of passing an initial seed to the
pseudo-random number generator, for each version. (These can be arbitrarily
typed in, or obtained by the methods discussed in
\textbf{\nameref{para:ControlSeed}} on page~\pageref*{para:ControlSeed}.)

\paragraph*{The problem.} For a multi-version \pkg{eqexam} document, where
each version samples from the same set of DB Test files, later versions of the exam
may have questions that are repeats of the ones posed in earlier versions of
the same source file. This may or may not be a problem if the different
versions are given to different classes at approximately the same time where
there is no opportunity for the details of the test to ``leak out'' from one
class to another.

\paragraph*{The solution.} The solution to this problem requires the introduction of several
new commands.
\bVerb\takeMeasure{\string\uniqueXDBChoicesOn}%
\setlength\eflength{\bxSize}%
\begin{minipage}{\bxSize}\kern0pt\def\1{\rlap{\sffamily\hskip\eflength\%\textsf{ to view the IDs of problems used}}}\relax
\def\2{\rlap{\sffamily\hskip\eflength\%\textsf{ preamble only (required)}}}\relax
\def\3{\rlap{\sffamily\hskip\eflength\%\textsf{ there is also \cs{uniqueXDBChoicesOff} (the default)}}}\relax
\begin{Verbatim}[frame=single,commandchars=!()]
!3\uniqueXDBChoicesOn
!2\InputUsedIDs
!1%\viewIDstrue
\end{Verbatim}
\end{minipage}\eVerb The \cs{uniqueXDBChoicesOn} command
brings some special code into the \cs{useRanTok} command; this special code
tries to select problems that have not already been selected by earlier
versions of the source document. Key to this selection process is that the
special code needs to know which questions were earlier selected; that is the
role of the \cs{InputUsedIDs}. This latter command uses the value declared in
\cs{numVersions}, so it must appear \emph{after} the declaration of the
number of versions.

When \cs{uniqueXDBChoicesOn} is expanded, as the source document is compiled,
an auxiliary file named \cs{jobname-ver\ameta{Ltr}.cut} is written. This
file keeps a record of the problem ID of the problems selected, believe it or
not. \cs{InputUsedIDs} inputs the appropriate CUT files:
\begingroup\tightsettings
\begin{itemize}
\item for version B, \cs{jobname-verA.cut} is input;
\item for version C, \cs{jobname-verA.cut}, \cs{jobname-verB.cut} are input;
\item for version D,  \cs{jobname-verA.cut}, \cs{jobname-verB.cut}, \cs{jobname-verC.cut} are input;
\item and so on.
\end{itemize}
\endgroup\noindent
In this way, when you compile version \ameta{Ltr}, the document inputs all
the information it needs about previous versions to make in informed choice.

\paragraph*{Workflow.} When compiling a multi-version \pkg{eqexam} document, do the following:
\begin{enumerate}
  \item Build each version in \emph{alphabetical order}, that is, compile with \cs{forVersion\darg{a}},
  then \cs{forVersion\darg{b}}, and so on.
  \item Rename PDF produced to reflect the version \ameta{Ltr}; eg, \verb~\jobname-verA.pdf~ or
  \verb!\jobname-sec02.pdf!, or whatever.\footnote{The demo file \texttt{ctrl-build.tex} shows
  how to build all versions and rename the final PDF files all from one controlling file.}
\end{enumerate}

%\paragraph*{Final comments.} \cs{uniqueXDBChoicesOn} does several things, it causes \cs{useRanTok}
%to use some special code to ensure cross-version uniqueness of choices, but
%it also initiates writing of the choice history to the auxiliary CUT files
%\cs{jobname-ver\ameta{Ltr}.cut}. By placing \cs{uniqueXDBChoicesOff} in
%the body of the document (outside a group, perhaps) that turns off the
%execution of the special code as well and writing to the CUT files. I believe
%you can later in the document turn this back on with \cs{uniqueXDBChoicesOn},
%though I have not done a lot of testing of this.

\newtopic\noindent
Now, I simply must get back to my retirement. \dps

\end{document}
